% =====================================================================
%  2.7 Integração entre as ferramentas
% =====================================================================
\subsection{Integração entre as ferramentas}\label{sec:integracao}

Cada ferramenta foi estudada separadamente, mas o objetivo final era entender como elas se combinam
em uma esteira DevOps (Figura~\ref{fig:integracao} e Tabela~\ref{tab:fases-ferramentas}).

\begin{figure}[H]
  \centering
  \caption{Integração entre as ferramentas no laboratório}\label{fig:integracao}
  \begin{tikzpicture}
    \tikzset{bloco/.style={caixa, text width=2.4cm, minimum height=1.2cm}}
    % Linha de cima: fluxo de CI/CD
    \node[caixacinza, text width=2.5cm, minimum height=1.2cm] (dev) at (0,0) {\textbf{Desenvolvedor}\\{\scriptsize edita e faz \textit{commit}}};
    \node[caixalaranja, text width=2.4cm, minimum height=1.2cm] (git) at (4.1,0) {\textbf{Git / GitHub}\\{\scriptsize versiona código e IaC}};
    \node[caixaescura, text width=2.4cm, minimum height=1.2cm] (jen) at (8.4,0) {Jenkins\\{\normalfont\scriptsize orquestra o \textit{pipeline}}};
    \node[bloco] (son) at (12.8,1.05) {\textbf{SonarQube}\\{\scriptsize qualidade e \textit{quality gate}}};
    \node[bloco] (doc) at (12.8,-1.05) {\textbf{Docker}\\{\scriptsize imagem e container}};
    % Linha de baixo: infraestrutura configurada
    \node[bloco] (ans) at (4.1,-3.1) {\textbf{Ansible}\\{\scriptsize \textit{control-node}}};
    \node[bloco] (app) at (8.4,-3.1) {\textbf{app01}\\{\scriptsize Java + Easy Notes}};
    \node[bloco] (db)  at (12.8,-3.1) {\textbf{db01}\\{\scriptsize MariaDB}};
    % Setas
    \draw[seta] (dev) -- node[rotulo, above=1pt]{\texttt{git push}} (git);
    \draw[seta] (git) -- node[rotulo, above=1pt]{\textit{poll}} node[rotulo, below=1pt]{checkout} (jen);
    \draw[seta] (jen.east) -- node[rotulo, above, sloped]{análise} (son.west);
    \draw[seta] (jen.east) -- node[rotulo, below, sloped]{\textit{build}/\textit{deploy}} (doc.west);
    \draw[seta, dashed] (git.south) -- node[rotulo, right=2pt]{\textit{playbooks}} (ans.north);
    \draw[seta] (ans) -- node[rotulo, above=1pt]{SSH} (app);
    \draw[seta] (ans.south) to[bend right=16] node[rotulo, above=1pt]{SSH} (db.south);
    \draw[seta, color=corverde] (app) -- node[rotulo, above=1pt]{JDBC 3306} (db);
    \draw[seta, dashed, color=corverde] (jen.south) to[out=-90, in=150]
         node[rotulo, pos=0.45, below left=0pt]{testes (JDBC)} (db.north west);
    \draw[seta, dashed, color=corverde] (doc.south) -- node[rotulo, right=2pt]{JDBC 3306} (db.north);
    % Base: Vagrant
    \coordinate (bl) at ($(dev.west |- ans.south)+(0,-1.65)$);
    \coordinate (br) at ($(db.east |- ans.south)+(0,-1.05)$);
    \draw[draw=corprimaria, fill=corprimaria!10, rounded corners=3pt] (bl) rectangle (br);
    \node[font=\scriptsize\bfseries, text=corprimaria] at ($(bl)!0.5!(br)$)
         {Vagrant + VMware: todas as VMs definidas em \texttt{Vagrantfile} (Infraestrutura como Código)};
  \end{tikzpicture}
  \fonte{Elaborado pelo autor (2026).}
\end{figure}

\begin{table}[H]
  \centering
  \caption{Etapas do ciclo DevOps e ferramentas utilizadas no laboratório}\label{tab:fases-ferramentas}
  \small
  \renewcommand{\arraystretch}{1.18}
  \begin{tabularx}{\linewidth}{>{\raggedright\arraybackslash}p{3.3cm} l >{\raggedright\arraybackslash}X}
    \toprule
    \textbf{Etapa} & \textbf{Ferramenta} & \textbf{Como foi usada} \\
    \midrule
    Codificar e versionar & Git e GitHub & Versionamento do código e de toda a configuração do laboratório. \\
    Provisionar a infraestrutura & Vagrant + VMware & Criação reprodutível das VMs a partir de arquivos de texto. \\
    Configurar servidores & Ansible & Preparação do sistema, MariaDB no \texttt{db01} e \textit{deploy} da aplicação no \texttt{app01}. \\
    \textit{Build} e testes & Maven (no Jenkins) & Compilação, empacotamento e testes a cada alteração. \\
    Qualidade e segurança & SonarQube & Análise estática e Quality Gate antes do empacotamento. \\
    Empacotar e \textit{deploy} & Docker + Jenkins & Imagem da aplicação e troca automática do container. \\
    Operar e monitorar & \texttt{systemd}, logs, \texttt{curl} & Verificação dos serviços e testes da API. \\
    \bottomrule
  \end{tabularx}
  \fonte{Elaborado pelo autor (2026).}
\end{table}

Na prática, o caminho de uma alteração fica assim: o desenvolvedor faz \cod|git push|; o Jenkins
detecta o novo \textit{commit}, compila e testa a aplicação contra o banco do \texttt{db01} --
configurado pelo Ansible em uma VM criada pelo Vagrant --; o SonarQube analisa o código e o Quality
Gate decide se ele pode seguir; se aprovado, o Docker empacota a aplicação e o Jenkins a coloca no ar.
Se algo falha, o \textit{pipeline} para e o \textit{feedback} chega rápido. O ponto central é que
\textbf{tudo é código} versionado no Git -- as máquinas, a configuração, a imagem e o próprio processo
de entrega --, o que torna o ambiente reprodutível e fácil de evoluir.
